\documentclass[11pt]{article}
\usepackage[margin=1in]{geometry}
\usepackage{amsmath,amssymb,amsthm}
\usepackage{booktabs}
\usepackage{hyperref}

\newtheorem{theorem}{Theorem}
\newtheorem{lemma}[theorem]{Lemma}
\newtheorem{corollary}[theorem]{Corollary}
\newtheorem{proposition}[theorem]{Proposition}
\theoremstyle{definition}
\newtheorem{definition}[theorem]{Definition}
\newtheorem{example}[theorem]{Example}
\newtheorem{remark}[theorem]{Remark}

\newcommand{\Con}{\operatorname{Con}}

\title{A disproof of conjecture \texttt{00000008540}\\
\large\upshape $\Con(L)$ can be Boolean while $L$ is not subdirectly
irreducible, already for the three-element chain}
\author{}
\date{2026-09-12}

\begin{document}
\maketitle

\begin{abstract}
Conjecture \texttt{00000008540} asserts that the congruence lattice of a
finite lattice is Boolean \emph{exactly when} the lattice is of subdirectly
irreducible type. We refute the ``only if'' direction with the smallest
possible witness: the three-element chain $C_3$. Its congruence lattice
$\Con(C_3)$ has four elements and is Boolean, yet $C_3$ is not subdirectly
irreducible --- it embeds subdirectly into $C_2 \times C_2$ via
$0 \mapsto (0,0)$, $1 \mapsto (0,1)$, $2 \mapsto (1,1)$, and neither
projection is an isomorphism. Equivalently, $\Con(C_3)$ has two atoms, whereas
subdirect irreducibility requires exactly one. No two-element lattice is a
counterexample, so $C_3$ is optimal in size.

We give a second witness, the four-element Boolean lattice $\mathbf{2}^2 = 2
\times 2$, which fails for the transparent reason that it is a direct product
of two nontrivial lattices. Both witnesses are verified against brute force
over \emph{all} partitions of the underlying set (the Bell numbers $B_3 = 5$
and $B_4 = 15$), using the full congruence compatibility condition. We also
record the count produced by the weaker, incorrect test ``each block is a
sublattice'', since it differs and a reviewer comparing implementations should
not be surprised by it. The ``at most $2^{n-1}$'' clause of the conjecture is
\emph{not} refuted here; we say so explicitly.
\end{abstract}

\section{The conjecture}

Quoted verbatim from \texttt{conjectures/00000008540.md}:

\begin{quote}
\textbf{Definition:} A congruence of a lattice is an order-preserving
equivalence. \textbf{Conjecture:} The congruence lattice of a finite lattice
is Boolean exactly when the lattice is of subdirectly irreducible type; the
number of congruences is at most $2^{n-1}$, and the bound is optimal.
(Boolean classification of congruence lattices)
\end{quote}

\subsection*{Interpretation used}

The conjecture's phrasing is not fully formal, so we state the reading we
refute and flag it explicitly.

\begin{enumerate}
\item ``of subdirectly irreducible type'' is read as \emph{subdirectly
      irreducible} in the standard sense: $L$ is nontrivial and admits no
      subdirect representation $L \le \prod_i L_i$ in which every projection
      is surjective but no projection is an isomorphism. For a finite lattice
      this is equivalent to: $\Con(L)$ has \emph{exactly one} atom.
\item ``Boolean'' is read as: $\Con(L)$ is a Boolean lattice, i.e.\
      $\Con(L) \cong \mathbf{2}^k$ for some $k \ge 0$.
\item ``exactly when'' is read as a genuine biconditional. Refuting either
      direction refutes the conjecture.
\end{enumerate}

Under this reading the biconditional fails, and the witness has three
elements. We note in passing that the conjecture is also non-trivial to read
because subdirect irreducibility is strictly stronger than direct
indecomposability for lattices; our witness $C_3$ is precisely the classical
illustration of that gap.

\section{Definitions}

\begin{definition}[congruence of a lattice]\label{def:congruence}
An equivalence relation $\theta$ on a lattice $L$ is a \emph{congruence} if
\[
x \mathrel{\theta} y \;\Longrightarrow\;
(x \wedge z) \mathrel{\theta} (y \wedge z)
\quad\text{and}\quad
(x \vee z) \mathrel{\theta} (y \vee z)
\]
for all $x, y, z \in L$. The set $\Con(L)$ of all congruences of $L$, ordered
by inclusion of relations, is a lattice: meets and joins are intersection and
transitive closure.
\end{definition}

\begin{remark}
The two clauses are not reducible to a single one, and neither is implied by
``each block is a sublattice''. This matters for computation: checking only
block-wise sublattice closure accepts spurious relations. See
Remark~\ref{rem:method}.
\end{remark}

\begin{definition}[atom; subdirectly irreducible]
An \emph{atom} of $\Con(L)$ is an element covering the diagonal
$\Delta = \{(x,x) : x \in L\}$, i.e.\ a minimal element strictly above
$\Delta$. A lattice $L$ with $|L| \ge 2$ is \emph{subdirectly irreducible}
(SI) if $\Con(L)$ has exactly one atom.
\end{definition}

\begin{definition}[the witnesses]\label{def:witnesses}
Let $C_n$ denote the $n$-element chain $0 < 1 < \dots < n-1$. Let
$\mathbf{2}^2$ denote the four-element Boolean lattice, with elements
$\{0, a, b, 1\}$ where $a$ and $b$ are incomparable, $0$ the bottom and $1$
the top.
\end{definition}

\section{The minimal counterexample: the chain \texorpdfstring{$C_3$}{C3}}

\begin{proposition}\label{prop:chain-con}
A partition of a chain $C_n$ is a congruence if and only if every block is an
interval. Consequently $|\Con(C_n)| = 2^{n-1}$, and in particular
\[
\Con(C_3) = \bigl\{\;\{0\}\!\mid\!\{1\}\!\mid\!\{2\},\quad
\{0,1\}\!\mid\!\{2\},\quad
\{0\}\!\mid\!\{1,2\},\quad
\{0,1,2\}\;\bigr\}
\]
has exactly four elements.
\end{proposition}

\begin{proof}
\emph{Blocks are intervals.} Suppose a block $B$ is not an interval, so there
are $x < z < y$ with $x, y \in B$ and $z \notin B$. Since $x \wedge z = x$ and
$y \wedge z = z$, compatibility applied to the related pair $(x,y)$ with
$z$ gives $x \mathrel{\theta} z$, contradicting $z \notin B$.

\emph{Interval partitions are congruences.} Let $\theta$ have interval blocks
and let $x \mathrel{\theta} y$, so $x$ and $y$ lie in a common interval block
$B$; assume $x \le y$, so $B$ contains every $u$ with $x \le u \le y$. For
arbitrary $z$: if $z \ge y$ then $x \wedge z = x$ and $y \wedge z = y$, both in
$B$; if $z \le x$ then both equal $z$; if $x < z < y$ then $x \wedge z = x$ and
$y \wedge z = z$, and $z \in B$ because $B$ is an interval. The three cases for
$\vee$ are symmetric. Hence $\theta$ is a congruence.

\emph{Counting.} A partition of the chain into intervals is determined by
choosing which of the $n-1$ gaps between consecutive elements are cut, so
there are $2^{n-1}$ of them. For $n = 3$ this gives $4$, and listing them
gives the four relations displayed.
\end{proof}

\begin{corollary}\label{cor:con3-boolean}
$\Con(C_3)$ is a Boolean lattice.
\end{corollary}

\begin{proof}
By Proposition~\ref{prop:chain-con}, $\Con(C_3)$ has four elements. In the
inclusion order, $\Delta = \{0\}\!\mid\!\{1\}\!\mid\!\{2\}$ is the bottom,
$\nabla = \{0,1,2\}$ is the top, and $\alpha = \{0,1\}\!\mid\!\{2\}$ and
$\beta = \{0\}\!\mid\!\{1,2\}$ are incomparable: their intersection is
$\Delta$ and their join is $\nabla$ (transitive closure merges all three
elements). A four-element distributive lattice with two incomparable middle
elements sharing a bottom and a top is $\mathbf{2}^2$. Hence
$\Con(C_3) \cong \mathbf{2}^2$, which is Boolean.
\end{proof}

\begin{proposition}\label{prop:c3-not-si}
$C_3$ is not subdirectly irreducible.
\end{proposition}

\begin{proof}
Define $\varphi : C_3 \to C_2 \times C_2$ by
\[
\varphi(0) = (0,0), \qquad \varphi(1) = (0,1), \qquad \varphi(2) = (1,1).
\]
This is injective, and it preserves the lattice operations: on a chain both
$\wedge$ and $\vee$ are the order-theoretic min and max, and the coordinatewise
operations of $C_2 \times C_2$ agree with those of $C_3$ on the image. For
instance $\varphi(0 \vee 1) = \varphi(1) = (0,1) = (0,0) \vee (0,1) =
\varphi(0) \vee \varphi(1)$, and the remaining cases are immediate from
$0 < 1 < 2$ and $(0,0) < (0,1) < (1,1)$.

Both coordinate projections $\pi_1, \pi_2$ are surjective onto $C_2$: the
first coordinates of the image are $\{0, 0, 1\} = C_2$, and the second are
$\{0, 1, 1\} = C_2$. So $C_3$ is a \emph{subdirect} product of two copies of
$C_2$. Neither projection is an isomorphism, since $|C_3| = 3 \neq 2 =
|C_2|$. Hence $C_3$ admits a subdirect representation in which no projection is
an isomorphism, i.e.\ $C_3$ is not subdirectly irreducible.
\end{proof}

\begin{corollary}\label{cor:c3-atoms}
$\Con(C_3)$ has exactly two atoms, namely $\alpha = \{0,1\}\!\mid\!\{2\}$ and
$\beta = \{0\}\!\mid\!\{1,2\}$.
\end{corollary}

\begin{proof}
By Proposition~\ref{prop:chain-con} the only congruences are
$\Delta, \alpha, \beta, \nabla$, and both $\alpha$ and $\beta$ cover $\Delta$
(the remaining element $\nabla$ is strictly above each of them, since
$\alpha \subsetneq \nabla$ and $\beta \subsetneq \nabla$). So there are exactly
two minimal elements strictly above $\Delta$.
\end{proof}

\begin{theorem}\label{thm:main}
Conjecture \texttt{00000008540} is false.
\end{theorem}

\begin{proof}
By Corollary~\ref{cor:con3-boolean}, $\Con(C_3)$ is Boolean. By
Proposition~\ref{prop:c3-not-si}, $C_3$ is not subdirectly irreducible.
Corollary~\ref{cor:c3-atoms} records the same obstruction in the standard
form: $\Con(C_3)$ has two atoms, while Definition~2.3 requires exactly one.

Thus $\Con(C_3)$ is Boolean while $C_3$ is not subdirectly irreducible,
refuting the ``only if'' half of the biconditional.
\end{proof}

\section{Minimality}

\begin{proposition}\label{prop:minimal}
No lattice with fewer than three elements is a counterexample; $C_3$ is
therefore the smallest witness.
\end{proposition}

\begin{proof}
A lattice with one element is trivial and we exclude it from the SI
condition (Definition~2.3 requires $|L| \ge 2$); its congruence lattice has
one element, which is Boolean, but the question of whether it is SI is a
boundary convention and does not affect the claim.

The only lattice with two elements is $C_2$. Its congruences are $\Delta$ and
$\nabla$, so $\Con(C_2)$ has two elements, exactly one atom, and $C_2$ is
subdirectly irreducible. Hence $C_2$ satisfies the conjecture and is not a
counterexample.
\end{proof}

\begin{remark}
The chain $C_3$ is directly indecomposable: a direct product of two nontrivial
lattices has at least four elements, and $C_3$ is a chain whereas such a
product never is one. So $C_3$ is directly indecomposable but not subdirectly
irreducible, which is why the ``only if'' direction of the conjecture is
vulnerable precisely here. This is the classical illustration that subdirect
irreducibility is strictly stronger than direct indecomposability.
\end{remark}

\section{A second counterexample}

The same failure occurs for the four-element Boolean lattice, by a more
transparent argument.

\begin{proposition}\label{prop:22}
$\mathbf{2}^2$ is not subdirectly irreducible, while $\Con(\mathbf{2}^2)$ is
Boolean.
\end{proposition}

\begin{proof}
\emph{Not SI.} $\mathbf{2}^2$ is by construction the direct product
$2 \times 2$ of two nontrivial lattices, so it is subdirectly reducible. Its
congruence lattice is $\Con(2) \times \Con(2) = 2 \times 2$, which has two
atoms; equivalently, the two projections $2 \times 2 \to 2$ are surjective and
neither is injective.

\emph{Boolean.} The congruences of $\mathbf{2}^2$ are exactly the kernels of
the lattice homomorphisms out of it:
\[
\Delta = \{0\}\!\mid\!\{a\}\!\mid\!\{b\}\!\mid\!\{1\},\quad
\alpha = \{0,a\}\!\mid\!\{b,1\},\quad
\beta = \{0,b\}\!\mid\!\{a,1\},\quad
\nabla = \{0,a,b,1\}.
\]
Each is visibly an equivalence, and each is a kernel of a homomorphism
($\Delta$ of the identity, $\nabla$ of the constant map, $\alpha$ of the
projection onto the $b$-coordinate, $\beta$ of the projection onto the
$a$-coordinate), hence a congruence. Conversely, if $\theta$ is a congruence
and $a \mathrel{\theta} b$ then $1 = a \vee b \mathrel{\theta} b \vee b = b$
and $0 = a \wedge b \mathrel{\theta} b$, so $0 \mathrel{\theta} 1$ and
$\theta = \nabla$. Otherwise $a, b$ are separated; if $0 \mathrel{\theta} a$
then $b = 0 \vee b \mathrel{\theta} a \vee b = 1$, giving
$\theta = \alpha$, and symmetrically $0 \mathrel{\theta} b$ gives
$\theta = \beta$; if neither holds then $\theta = \Delta$. So there are exactly
four, forming the diamond, so $\Con(\mathbf{2}^2) \cong \mathbf{2}^2$ is
Boolean.
\end{proof}

\section{The counting clause is not refuted}

\begin{remark}[scope of the refutation]\label{rem:counting}
The conjecture also asserts $|\Con(L)| \le 2^{n-1}$ for $n = |L|$, with the
bound optimal. Our counterexamples do not touch this clause:
\[
|\Con(C_3)| = 4 = 2^{3-1}, \qquad |\Con(\mathbf{2}^2)| = 4 \le 2^{4-1} = 8.
\]
In fact the bound is attained by every chain, by
Proposition~\ref{prop:chain-con}, which is consistent with the conjecture's
own claim that the bound is optimal.

We state this explicitly so that the scope of the refutation is unambiguous:
what fails is the \emph{classification} (the biconditional), not the
\emph{counting} bound. A reader should not infer from ``the conjecture is
false'' that every clause is false.
\end{remark}

\section{Verification}

Both enumerations were checked by brute force over \emph{all} partitions of
the underlying set, testing the full compatibility condition of
Definition~2.1, i.e.\ for every related pair $(x,y)$ and every $z$, both
$(x \wedge z, y \wedge z)$ and $(x \vee z, y \vee z)$ must be related.

\begin{table}[h]
\centering
\begin{tabular}{lrrrrrr}
\toprule
$L$ & $|L|$ & partitions tested & wrong test & full condition & atoms & SI\\
\midrule
$C_2$ & 2 & 2  & 2  & 2 & 1 & yes\\
$C_3$ & 3 & 5  & 5  & \textbf{4} & \textbf{2} & \textbf{no}\\
$\mathbf{2}^2$ & 4 & 15 & 11 & \textbf{4} & \textbf{2} & \textbf{no}\\
\bottomrule
\end{tabular}
\caption{Brute-force counts. ``partitions tested'' is the Bell number $B_n$
($B_2 = 2$, $B_3 = 5$, $B_4 = 15$). ``wrong test'' is the count obtained by
the incorrect condition ``each block is a sublattice''; it overcounts, and the
discrepancy is recorded deliberately. Note that $B_4 = 15$, not $256$: $256$
is neither the number of partitions of a four-element set nor the number of
equivalence relations on it (the two counts coincide at $15$, since
partitions and equivalence relations are in bijection).}
\end{table}

\begin{remark}[the wrong test]\label{rem:method}
For $C_3$ and $\mathbf{2}^2$ the weak test ``each block is a sublattice''
accepts spurious relations: it constrains $x \wedge y$ and $x \vee y$ for
related $x, y$ but says nothing about $x \wedge z$ for an unrelated $z$. It
reports $5$ instead of $4$ for $C_3$ --- the spurious relation being
$\{0,2\}\!\mid\!\{1\}$, whose block $\{0,2\}$ is indeed a sublattice of $C_3$
but which is not a congruence, as the proof of
Proposition~\ref{prop:chain-con} shows --- and $11$ instead of $4$ for
$\mathbf{2}^2$. We record the wrong intermediate values so that a reviewer
comparing implementations is not surprised by them.
\end{remark}

\begin{remark}[$C_3$ is not the only chain that works]\label{rem:chains}
The argument of Theorem~\ref{thm:main} generalises: for every $n \ge 3$ the
chain $C_n$ is a counterexample. By Proposition~\ref{prop:chain-con},
$\Con(C_n)$ has $2^{n-1}$ elements and is Boolean, while its atoms are the
$n-1$ interval partitions obtained by merging a single adjacent pair, so it
has $n-1 \ge 2$ atoms and $C_n$ is not SI. We present $C_3$ because
Proposition~\ref{prop:minimal} shows it is the smallest witness, and because
the failure is easiest to see there; $\mathbf{2}^2$ is included because its
non-SI-ness is immediate from being a direct product. Both are given so that
the reader has a witness whose non-SI-ness is structural ($\mathbf{2}^2$) and
one whose size is minimal ($C_3$).
\end{remark}

\subsection*{Formalisation}

The combinatorial core of the refutation is formalised in Lean~4 in the
accompanying \texttt{lean4/} project: that $C_3$ has exactly the four
congruences listed above, that they are order isomorphic to
$\mathbf{2} \times \mathbf{2}$, and that exactly two of them are atoms. It
builds with \texttt{lake build} with no \texttt{sorry}, using \emph{core Lean
only} --- no Mathlib --- and every theorem depends solely on the standard
axioms \texttt{propext} and \texttt{Quot.sound}. The two statements that carry
the refutation are

\begin{quote}\ttfamily
theorem con\_complete (f : Fin 3 → Fin 3 → Bool) (h : IsCongruence f) :\\
\hspace*{1.5em}f = relBot $\vee$ f = relAlpha $\vee$ f = relBeta $\vee$ f = relTop\\[0.5em]
theorem atom\_only\_two (f : Fin 3 → Fin 3 → Bool) (h : IsAtom f) :\\
\hspace*{1.5em}f = relAlpha $\vee$ f = relBeta
\end{quote}

\noindent Booleanity is formalised separately, as the triple
\texttt{code\_inj}, \texttt{code\_surj}, \texttt{code\_le}: the map
$f \mapsto (f\,0\,1, f\,1\,2)$ is a bijection from $\Con(C_3)$ onto
$\mathtt{Bool} \times \mathtt{Bool}$ that preserves and reflects the order, and
$\mathtt{Bool} \times \mathtt{Bool}$ with the componentwise order is the
four-element Boolean lattice $\mathbf{2}^2$.

Congruences are represented as \texttt{Bool}-valued relations, so that every
clause of Definition~\ref{def:congruence} is a decidable proposition and the
concrete instances close by \texttt{decide}. The mathematical content is
carried by two lemmas: \texttt{f02\_eq} (on a chain, $0 \sim 2$ holds exactly
when both $0 \sim 1$ and $1 \sim 2$ do) and \texttt{determined} (a congruence on
a chain is determined by the pair $(f\,0\,1, f\,1\,2)$), after which
\texttt{con\_complete} is a four-case split.

The formalisation covers the enumeration and the atom count --- precisely the
part a reviewer would otherwise have to check by hand. It does \emph{not}
formalise the surrounding universal algebra: ``finite lattice'', ``subdirect
product'' and ``subdirectly irreducible'' are used as standard notions rather
than reproduced in Lean. It also does not formalise the counting clause, which
is not refuted.

\section{Reproducibility}

Pure Python~3, standard library only:

\begin{quote}\ttfamily
\$ python3 reproduce.py
\end{quote}

The script enumerates all partitions of $\{0,\dots,n-1\}$ as restricted growth
strings, filters by the full congruence condition, builds the congruence
lattice, counts atoms, and tests Booleanity as ``$|\Con(L)| = 2^{k}$ with every
element a join of atoms'', where $k$ is the atom count. It reports $C_3$ and
$\mathbf{2}^2$ as counterexamples, confirms that $C_2$ is not one, and prints
both the correct and the incorrect counts. No step of the proof depends on the
computation.

\end{document}
